% GENERATED FILE. Edit canonical lessons, then run npm run generate:books.
% canonical-source-hash: fnv1a64:b0f55af7bc7c6708
% canonical-lessons: GE-C211-aengstlich, GE-C211-anstrengend, GE-C211-aufgeregt, GE-C211-ausgezeichnet, GE-C211-berufstaetig

\chapter{Anxious, Tiring, Excited, Excellent, Working}
\label{ch:ge-anxious-tiring-excited-excellent-working}
\hlchaptermodality{211}

\begin{chapteropening}
\textbf{By the end of this chapter:} \emph{I can say anxious, tiring, excited, excellent and working.}

Complete the last lesson of chapter 211: I can say anxious, tiring, excited, excellent and working.
\end{chapteropening}

\section[ängstlich]{ängstlich — anxious, timid}

\label{lesson:GE-C211-aengstlich}



\begin{quote}
Before the new one: say the German for a butcher's, then the German for a department store.
\end{quote}

\subsection*{You'll want to know: ängstlich}
\textbf{ängstlich} — \textquotedblleft{}anxious, timid\textquotedblright{}.

\begin{cousinweb}
\textbf{ängstlich} is \emph{Angst}, \textquotedblleft{}fear\textquotedblright{}, + \emph{-lich}: fearful. \emph{Angst} $\leftarrow$ Proto-Germanic \emph{*angustiz}, from a word for \textquotedblleft{}narrow\textquotedblright{}. (EtymDB 2.1)
\end{cousinweb}

\begin{grammarlens}[title={What you've built}]
One more word for this chapter.
\end{grammarlens}

\subsection*{Your turn}
\begin{itemize}
\raggedright
  \item \emph{Say it:} \emph{ängstlich}
  \item \emph{Say it:} \emph{ängstlich}, once more
  \item \emph{Say it:} say \emph{ängstlich}
  \item \emph{Recall:} say the German for a butcher's, then the German for a department store, then say \emph{ängstlich} again
\end{itemize}

\begin{tcolorbox}[breakable,colback=teal!4,colframe=teal!35!black,title={Before you move on}]
What does ängstlich mean? (\textquotedblleft{}Anxious, timid\textquotedblright{}.) Say it once more.
\end{tcolorbox}
\section[anstrengend]{anstrengend — tiring, demanding}

\label{lesson:GE-C211-anstrengend}



\begin{quote}
Before the new one: say the German for a department store, then the German for anxious.
\end{quote}

\subsection*{You'll want to know: anstrengend}
\textbf{anstrengend} — \textquotedblleft{}tiring, demanding\textquotedblright{}.

\begin{cousinweb}
\textbf{anstrengend} is the present participle of the verb \textquotedblleft{}to strain\textquotedblright{}, built on the word for \textquotedblleft{}tight, strict\textquotedblright{}, Old High German \emph{strengi}, which EtymDB 2.1 lists beside English \emph{strong}.
\end{cousinweb}

\begin{grammarlens}[title={What you've built}]
One more word for this chapter.
\end{grammarlens}

\subsection*{Your turn}
\begin{itemize}
\raggedright
  \item \emph{Say it:} \emph{anstrengend}
  \item \emph{Say it:} \emph{anstrengend}, once more
  \item \emph{Say it:} say \emph{anstrengend}
  \item \emph{Recall:} say the German for a department store, then the German for anxious, then say \emph{anstrengend} again
\end{itemize}

\begin{tcolorbox}[breakable,colback=teal!4,colframe=teal!35!black,title={Before you move on}]
What does anstrengend mean? (\textquotedblleft{}Tiring, demanding\textquotedblright{}.) Say it once more.
\end{tcolorbox}
\section[aufgeregt]{aufgeregt — excited, nervous}

\label{lesson:GE-C211-aufgeregt}



\begin{quote}
Before the new one: say the German for anxious, then the German for tiring.
\end{quote}

\subsection*{You'll want to know: aufgeregt}
\textbf{aufgeregt} — \textquotedblleft{}excited, nervous\textquotedblright{}.

\begin{cousinweb}
\textbf{aufgeregt} is the past participle of \emph{auf} + an old verb meaning \textquotedblleft{}to stir, to move\textquotedblright{}: stirred up.
\end{cousinweb}

\begin{grammarlens}[title={What you've built}]
One more word for this chapter.
\end{grammarlens}

\subsection*{Your turn}
\begin{itemize}
\raggedright
  \item \emph{Say it:} \emph{aufgeregt}
  \item \emph{Say it:} \emph{aufgeregt}, once more
  \item \emph{Say it:} say \emph{aufgeregt}
  \item \emph{Recall:} say the German for anxious, then the German for tiring, then say \emph{aufgeregt} again
\end{itemize}

\begin{tcolorbox}[breakable,colback=teal!4,colframe=teal!35!black,title={Before you move on}]
What does aufgeregt mean? (\textquotedblleft{}Excited, nervous\textquotedblright{}.) Say it once more.
\end{tcolorbox}
\section[ausgezeichnet]{ausgezeichnet — excellent}

\label{lesson:GE-C211-ausgezeichnet}



\begin{quote}
Before the new one: say the German for tiring, then the German for excited.
\end{quote}

\subsection*{You'll want to know: ausgezeichnet}
\textbf{ausgezeichnet} — \textquotedblleft{}excellent\textquotedblright{}.

\begin{cousinweb}
\textbf{ausgezeichnet} is the past participle of \emph{aus-} + \emph{zeichnen}, \textquotedblleft{}to mark\textquotedblright{}: marked out from the rest. \emph{zeichnen} is built on the word for \textquotedblleft{}sign\textquotedblright{}, English \emph{token}. (EtymDB 2.1)
\end{cousinweb}

\begin{grammarlens}[title={What you've built}]
One more word for this chapter.
\end{grammarlens}

\subsection*{Your turn}
\begin{itemize}
\raggedright
  \item \emph{Say it:} \emph{ausgezeichnet}
  \item \emph{Say it:} \emph{ausgezeichnet}, once more
  \item \emph{Say it:} say \emph{ausgezeichnet}
  \item \emph{Recall:} say the German for tiring, then the German for excited, then say \emph{ausgezeichnet} again
\end{itemize}

\begin{tcolorbox}[breakable,colback=teal!4,colframe=teal!35!black,title={Before you move on}]
What does ausgezeichnet mean? (\textquotedblleft{}Excellent\textquotedblright{}.) Say it once more.
\end{tcolorbox}
\section[berufstätig]{berufstätig — working, employed}

\label{lesson:GE-C211-berufstaetig}



\begin{quote}
Before the new one: say the German for excited, then the German for excellent.
\end{quote}

\subsection*{You'll want to know: berufstätig}
\textbf{berufstätig} — \textquotedblleft{}working, employed\textquotedblright{}.

\begin{cousinweb}
\textbf{berufstätig} is \emph{Beruf}, \textquotedblleft{}profession\textquotedblright{}, + \emph{-s-} + the word for \textquotedblleft{}active\textquotedblright{}, built on the word for \textquotedblleft{}a deed\textquotedblright{}, Old High German \emph{tāt}, English \emph{deed} (EtymDB 2.1).
\end{cousinweb}

\begin{grammarlens}[title={What you've built}]
One more word for this chapter.
\end{grammarlens}

\subsection*{Your turn}
\begin{itemize}
\raggedright
  \item \emph{Say it:} \emph{berufstätig}
  \item \emph{Say it:} \emph{berufstätig}, once more
  \item \emph{Say it:} say \emph{berufstätig}
  \item \emph{Recall:} say the German for excited, then the German for excellent, then say \emph{berufstätig} again
\end{itemize}

\begin{tcolorbox}[breakable,colback=teal!4,colframe=teal!35!black,title={Before you move on}]
What does berufstätig mean? (\textquotedblleft{}Working, employed\textquotedblright{}.) Say it once more.
\end{tcolorbox}
